% =============================================================================
% Project : THE_PHILOSOPHY_OF_VIABILITY
% Book    : Book I — Ontology of Viability
% Book ID : VS.B01
%
% File    : 3_BACKMATTER/02_logical_entity_reference_map.tex
%
% Component:
%   Logical Entity Reference Map
%
% Purpose:
%   Establishes the centralized documentary reference framework
%   for identifying, locating, classifying, and cross-referencing
%   the structural, logical, philosophical, and asset-based
%   entities belonging to Book I.
%
% Scope:
%   Project and book identification;
%   structural hierarchy and documentary locations;
%   independent asset identification;
%   logical and philosophical register categories;
%   epistemic and justificatory classifications;
%   reference relations and identifier governance.
%
% Dependencies:
%   Standard LaTeX commands;
%   array, tabularx, booktabs, xurl, hyperref, seqsplit.
%
% =============================================================================


% =============================================================================
% LOGICAL ENTITY REFERENCE MAP
% =============================================================================

\cleardoublepage
\chapter*{Logical Entity Reference Map}
\phantomsection
\addcontentsline{toc}{chapter}{Logical Entity Reference Map}

\markboth
    {Logical Entity Reference Map}
    {Logical Entity Reference Map}


% =============================================================================
% 01 — PURPOSE AND SCOPE
% =============================================================================

\section*{Purpose and Scope}

The Logical Entity Reference Map establishes a centralized
documentary framework for identifying and locating the principal
entities involved in the formal organization of Book I.

Its primary function is referential. It coordinates the identifiers
used for the project, the book, its structural divisions, its
independently maintained assets, and its logical and philosophical
registers.

The reference system makes it possible to locate a designated
entity through a consistent identifier while preserving the
distinction between the identifier, the entity to which it refers,
and the substantive content of that entity.

The scope of this map comprises six principal areas:

\begin{enumerate}
    \item Project and book identification.
    \item The structural organization of the book.
    \item The identification of chapters, sections, subsections,
          and formally designated textual units.
    \item The identification and centralized management of
          independent tables, figures, diagrams, and illustrations.
    \item The functional classification and identification of
          entries in logical and philosophical registers.
    \item The distinction between documentary location, logical
          classification, and epistemic or justificatory status.
\end{enumerate}

The map is a component of the documentary infrastructure of
the project. It does not establish new ontological principles,
introduce substantive philosophical categories, or determine
the truth, validity, or justificatory adequacy of registered
philosophical claims.

Its classifications organize the documentation of the
investigation; they do not replace the arguments through which
the philosophical commitments of the book are examined.


% =============================================================================
% 02 — PROJECT AND BOOK IDENTIFICATION
% =============================================================================

\section*{Project and Book Identification}

The project and book identifiers establish the primary
documentary reference levels of the philosophical library.

\begin{center}
\small
\setlength{\tabcolsep}{4pt}
\begin{tabularx}{\textwidth}{
    >{\raggedright\arraybackslash}p{0.13\textwidth}
    >{\raggedright\arraybackslash}p{0.20\textwidth}
    >{\raggedright\arraybackslash}X
}
\toprule
\textbf{Entity} &
\textbf{Identifier} &
\textbf{Reference Function}
\\
\midrule

Project
&
\texttt{VS}
&
Identifies the project entitled
\textit{The Philosophy of Viability}.
\\

Book I
&
\texttt{VS.B01}
&
Identifies the first book,
\textit{Ontology of Viability}.
\\

Part I
&
\texttt{VS.B01.P01}
&
Identifies the first major part of Book I.
\\

Chapter 01
&
\texttt{VS.B01.P01.C01}
&
Identifies Chapter 01, \textit{Identity}.
\\
\bottomrule
\end{tabularx}
\end{center}

These identifiers establish documentary and structural
references. Their hierarchical arrangement specifies
organizational containment within the book; it does not,
by itself, establish ontological dependence, philosophical
importance, or logical priority among the subjects discussed.

The identification of a documentary unit must therefore
be distinguished from the philosophical determination of
the entity, concept, or claim examined within that unit.


% =============================================================================
% 03 — STRUCTURAL HIERARCHY
% =============================================================================

\section*{Structural Hierarchy}

The structural identifiers of Book I follow a hierarchical
notation in which successive components specify increasingly
localized positions within the documentary organization.

A representative full-length structural identifier is:

\begin{center}
\small
\setlength{\tabcolsep}{4pt}
\renewcommand{\arraystretch}{1.15}
\resizebox{\textwidth}{!}{%
\begin{tabular}{@{}p{2.2cm}p{3.3cm}p{9.2cm}@{}}
\toprule
\textbf{Component} &
\textbf{Structural Level} &
\textbf{Reference Function}
\\
\midrule
\texttt{VS}   & Project    & Identifies \textit{The Philosophy of Viability}. \\
\texttt{B01}  & Book       & Identifies Book I. \\
\texttt{P01}  & Part       & Identifies Part I. \\
\texttt{C01}  & Chapter    & Identifies Chapter 01 within the designated part. \\
\texttt{S01}  & Section    & Identifies Section 01 within the designated chapter. \\
\texttt{SS01} & Subsection & Identifies Subsection 01 within its designated section. \\
\texttt{P001} & Paragraph  & Identifies a formally numbered paragraph within its designated structural location. \\
\bottomrule
\end{tabular}%
}
\end{center}

The complete identifier provides a documentary address
for a designated textual unit. The presence of a component
indicates that the corresponding structural level has been
specified in the identifier; it does not imply that every
textual unit must necessarily possess every possible level
of subdivision.

In particular, paragraph-level identification is required
only where the project's referencing practices call for
that degree of precision.

The structural hierarchy must not be interpreted as a
hierarchy of ontological dependence, philosophical authority,
or logical priority. It expresses the organization of the
document, not the structure of reality investigated by it.


% =============================================================================
% 04 — STRUCTURAL ENTITIES
% =============================================================================

\section*{Structural Entities}

The principal structural entities are distinguished
according to their respective documentary functions.

\begin{center}
\small
\setlength{\tabcolsep}{4pt}
\begin{tabularx}{\textwidth}{
    >{\raggedright\arraybackslash}p{0.20\textwidth}
    >{\raggedright\arraybackslash}p{0.23\textwidth}
    >{\raggedright\arraybackslash}X
}
\toprule
\textbf{Entity Type} &
\textbf{Identifier Component} &
\textbf{Documentary Function}
\\
\midrule

Book
&
\texttt{B}
&
Identifies an individual book within the project.
\\

Part
&
\texttt{P}
&
Identifies a major organizational division of a book.
\\

Chapter
&
\texttt{C}
&
Identifies a chapter within a designated part.
\\

Section
&
\texttt{S}
&
Identifies a principal section within a chapter.
\\

Subsection
&
\texttt{SS}
&
Identifies a subordinate textual division within a section.
\\

Paragraph
&
\texttt{P}
&
Identifies a formally numbered textual unit.
\\
\bottomrule
\end{tabularx}
\end{center}

The component \texttt{P} is used at two distinct structural
levels: Part and Paragraph. Its interpretation is therefore
context-dependent.

The distinction is determined by its position within the
complete identifier and by the form of its numerical suffix:

\begin{itemize}
    \item \texttt{P01} identifies a Part.
    \item \texttt{P001} identifies a Paragraph.
\end{itemize}

Consequently, the component must never be interpreted
in isolation when resolving a structural reference.

The complete identifier, rather than any single component,
determines the documentary location being designated.


% =============================================================================
% 05 — INDEPENDENT ASSET IDENTIFICATION
% =============================================================================

\section*{Independent Asset Identification}

Tables, figures, diagrams, and illustrations are maintained
as independently identified project assets.

Unlike ordinary textual passages, these assets possess
their own files and identifiers. Their physical storage
location is centralized, while their intended placement
within the book is specified through their structural
reference components.

An asset identifier combines:

\begin{enumerate}
    \item The project and book identifiers.
    \item The structural location with which the asset is
          associated.
    \item The asset-type code.
    \item A sequential number within the relevant
          identification scheme.
\end{enumerate}

A representative asset identifier is:

\begin{center}
\small
\path{VS.B01.P01.C01.S01.SS01.P001.TL01}
\end{center}

In this example, the final component identifies the asset
as the first table associated with the designated textual
location.

The principal asset codes are:

\begin{center}
\small
\setlength{\tabcolsep}{4pt}
\begin{tabularx}{\textwidth}{
    >{\raggedright\arraybackslash}p{0.13\textwidth}
    >{\raggedright\arraybackslash}p{0.20\textwidth}
    >{\raggedright\arraybackslash}X
}
\toprule
\textbf{Code} &
\textbf{Asset Type} &
\textbf{Reference Function}
\\
\midrule

\texttt{TL}
&
Table
&
Presents structured comparisons, classifications,
distinctions, or tabular information.
\\

\texttt{FG}
&
Figure
&
Presents a visual representation of a configuration,
structure, or other relevant subject.
\\

\texttt{DG}
&
Diagram
&
Represents conceptual, structural, relational,
or process-related configurations.
\\

\texttt{IL}
&
Illustration
&
Provides a visual or illustrative representation
intended to clarify a concept, relation, or example.
\\
\bottomrule
\end{tabularx}
\end{center}

The asset-type codes distinguish four categories of
independently maintained visual or structured resources.
They are identification categories rather than claims
about the philosophical function or evidential status
of the assets.

\subsection*{Centralized Asset Directory}

All independent assets belonging to Book I are maintained
within the central asset directory:

\begin{verbatim}
4_ASSETS/
    tables/
    figures/
    diagrams/
    illustrations/
\end{verbatim}

The four directories correspond to the asset-type codes:

\begin{itemize}
    \item \texttt{tables/} for \texttt{TL} assets.
    \item \texttt{figures/} for \texttt{FG} assets.
    \item \texttt{diagrams/} for \texttt{DG} assets.
    \item \texttt{illustrations/} for \texttt{IL} assets.
\end{itemize}

The physical location of an asset within
\texttt{4\_ASSETS/} must be distinguished from its
documentary association with a chapter, section,
subsection, or paragraph.

Centralized storage does not make an asset an ordinary
textual component of the section in which it is cited.
The asset remains an independently maintained project
entity whose intended placement is established through
references in the relevant textual file.

The structural components included in an asset identifier
must correspond to the asset's actual designated placement.
A paragraph-level component is used when the asset is
associated with a formally identified paragraph; it must
not be added merely to make an identifier appear more
complete.

Asset identifiers must remain consistent across:

\begin{enumerate}
    \item The asset's filename.
    \item The identifier recorded within the asset file,
          where applicable.
    \item The reference inserted into the relevant
          chapter or section.
    \item The corresponding project registry.
\end{enumerate}

Any relocation or change in the intended structural
association of an asset requires review of its references
and registry entry.


% =============================================================================
% 06 — LOGICAL AND PHILOSOPHICAL REGISTER IDENTIFICATION
% =============================================================================

\section*{Logical and Philosophical Register Identification}

The project maintains separate registers for entries
classified according to their logical or philosophical
function.

The established functional codes are:

\begin{center}
\small
\setlength{\tabcolsep}{4pt}
\begin{tabularx}{\textwidth}{
    >{\raggedright\arraybackslash}p{0.13\textwidth}
    >{\raggedright\arraybackslash}p{0.25\textwidth}
    >{\raggedright\arraybackslash}X
}
\toprule
\textbf{Code} &
\textbf{Register Category} &
\textbf{Reference Function}
\\
\midrule

\texttt{AX}
&
Axiom
&
Identifies an entry classified as an axiom.
\\

\texttt{DF}
&
Definition
&
Identifies an entry classified as a definition.
\\

\texttt{PR}
&
Proposition
&
Identifies an entry classified as a proposition.
\\

\texttt{TH}
&
Thesis
&
Identifies an entry classified as a thesis.
\\

\texttt{HY}
&
Hypothesis
&
Identifies an entry classified as a hypothesis.
\\

\texttt{AR}
&
Argument
&
Identifies an entry classified as an argument.
\\

\texttt{PF}
&
Philosophical Foundation
&
Identifies an entry classified as a philosophical
foundation.
\\
\bottomrule
\end{tabularx}
\end{center}

These codes identify functional or logical categories.
They must not be interpreted as a unified scale of
philosophical authority, evidential strength, or
epistemic justification.

The categories differ in their logical functions:

\begin{itemize}
    \item An \textbf{axiom} occupies a designated foundational
          role within a formal or conceptual system.
    \item A \textbf{definition} specifies the adopted or
          stipulated use of a term.
    \item A \textbf{proposition} expresses a claim presented
          for examination or assessment.
    \item A \textbf{thesis} expresses a substantive position
          advanced and defended within the philosophical
          investigation.
    \item A \textbf{hypothesis} identifies a proposed claim
          whose status remains subject to examination.
    \item An \textbf{argument} is a form of reasoning through
          which premises are related to a conclusion.
    \item A \textbf{philosophical foundation} identifies a
          designated foundational commitment or basis of
          the philosophical framework.
\end{itemize}

These distinctions concern the functional classification
of entries. They do not independently determine whether
an entry is true, adequately justified, formally valid,
or philosophically established.

In particular, the classification of a claim as an axiom
does not by itself establish its ontological truth.
Likewise, the classification of an entry as a definition
does not establish the ontological validity of the
conceptual commitments associated with that definition.

The functional category and the justificatory status
of an entry must therefore be recorded and evaluated
as distinct dimensions.


% =============================================================================
% 07 — EPISTEMIC AND JUSTIFICATORY STATUS
% =============================================================================

\section*{Epistemic and Justificatory Status}

The project distinguishes the logical or functional form
of a registered entry from its epistemic and justificatory
status.

The following eight categories identify the principal
justificatory positions used in the evaluation of
philosophical claims:

\begin{enumerate}
    \item \textbf{Foundationally Adopted Commitments.}
          Commitments adopted as part of the foundational
          framework of the investigation.

    \item \textbf{Stipulated or Defined Claims.}
          Claims or formulations whose terms or use have
          been explicitly stipulated or defined.

    \item \textbf{Argued or Derived Claims.}
          Claims supported through explicit reasoning
          or derived from previously established commitments.

    \item \textbf{Conditionally Supported Claims.}
          Claims supported under specified conditions,
          assumptions, or restrictions.

    \item \textbf{Provisional Constructions.}
          Conceptual or structural constructions retained
          provisionally and subject to further examination.

    \item \textbf{Hypotheses.}
          Proposed claims whose adequacy or truth remains
          to be investigated.

    \item \textbf{Unresolved or Open Problems.}
          Questions or claims for which no adequate
          determination has yet been established.

    \item \textbf{Rejected or Superseded Claims.}
          Claims that have been rejected, replaced,
          or rendered obsolete by subsequent revisions
          or arguments.
\end{enumerate}

These categories describe the justificatory position
assigned to an entry. They do not specify its documentary
location or logical form.

The two dimensions must remain independent.

For example, identifying an entry as a proposition does
not determine whether it has been demonstrated, conditionally
supported, provisionally adopted, or left unresolved.

Similarly, classifying an entry as an axiom does not
automatically determine whether it is a foundationally
adopted commitment or whether its philosophical justification
has been independently examined.

The same functional category may therefore occur with
different justificatory statuses, provided that the
classification is appropriate to the particular entry.

The distinction is methodological and classificatory:
it prevents the logical form of a statement from being
mistaken for the degree or kind of justification
available for it.


% =============================================================================
% 08 — REFERENCE RELATIONS AND CROSS-REFERENCING
% =============================================================================

\section*{Reference Relations and Cross-Referencing}

A reference identifier provides a means of locating
a designated entity within the documentary organization
of the project.

The use of references is governed by the following
distinctions:

\begin{enumerate}
    \item \textbf{Structural reference.}
          Identifies the documentary location of a book,
          part, chapter, section, subsection, or textual unit.

    \item \textbf{Asset reference.}
          Identifies an independently maintained asset
          and its designated structural association.

    \item \textbf{Functional reference.}
          Identifies the logical or philosophical category
          assigned to a registered entry.

    \item \textbf{Justificatory reference.}
          Identifies the epistemic or argumentative status
          assigned to a claim or entry.

    \item \textbf{Cross-reference.}
          Establishes a navigational or documentary
          connection between two or more designated entities.
\end{enumerate}

These reference relations must not be treated as
interchangeable.

A reference to a chapter does not identify an individual
paragraph. An asset identifier does not replace the
identifier of the section in which the asset is cited.
A functional register code does not specify the structural
location of the registered entry. A cross-reference
does not, merely by connecting two entities, establish
a relation of logical dependence or ontological constitution
between them.

References must be sufficiently precise to identify
their intended targets. Where several entities are
associated with the same subject, the reference must
make clear which entity is being designated.

Where an entity is relocated, renamed, or structurally
reorganized, its references must be reviewed in accordance
with the project's identifier-preservation policy.

Existing identifiers must not be changed merely to
improve naming consistency when such a change would
unnecessarily invalidate established references.


% =============================================================================
% 09 — IDENTIFIER GOVERNANCE
% =============================================================================

\section*{Identifier Governance}

The creation, assignment, modification, and maintenance
of identifiers must follow consistent project-wide rules.

The following requirements govern identifier management:

\begin{enumerate}
    \item Every formally registered entity must receive
          an identifier appropriate to its entity type
          and reference namespace.

    \item Identifiers must conform to the established
          structural and naming conventions.

    \item Sequential numbering must be checked against
          the relevant existing registry before a new
          identifier is assigned.

    \item The same identifier must not be assigned to
          two distinct entities within the same
          identifier namespace.

    \item Existing identifiers must be preserved unless
          a documented structural or administrative
          revision requires modification.

    \item Relocations, renamings, and structural revisions
          must be reflected in the relevant references
          and registries.

    \item An identifier must not be treated as evidence
          of the truth, validity, philosophical importance,
          or justificatory authority of its referent.

    \item The retirement or replacement of an identifier
          must be documented so that historical references
          can be interpreted without ambiguity.
\end{enumerate}

Identifier uniqueness is a property of the maintained
project inventory, not a property that can be inferred
from the visual form of an identifier.

Accordingly, global uniqueness must be established
through verification against the complete project
inventory, including the relevant structural files,
asset directories, and registers.

The existence of a syntactically well-formed identifier
does not establish that the identifier has already been
registered, that it is unique, or that its assigned
referent is correctly located.


% =============================================================================
% 10 — LIMITS OF THE REFERENCE SYSTEM
% =============================================================================

\section*{Limits of the Reference System}

The reference system is a documentary and organizational
instrument. Its identifiers and classifications must
not be confused with the ontological structure of the
subject matter investigated in Book I.

The following limitations are therefore fundamental:

\begin{itemize}
    \item \textbf{Structural order does not establish
          ontological priority.}
          The order in which entities appear in the book
          does not determine their order of ontological
          dependence.

    \item \textbf{Numerical order does not establish
          philosophical importance.}
          A lower sequential number does not confer
          greater conceptual or philosophical significance.

    \item \textbf{Functional classification does not
          establish justificatory validity.}
          Assigning a logical category to an entry does
          not demonstrate the adequacy of that entry.

    \item \textbf{Documentary location does not establish
          conceptual dependence.}
          The placement of an asset or discussion within
          a section does not determine the ontological
          or logical relations of its subject matter.

    \item \textbf{Cross-reference does not establish
          necessary relation.}
          A documentary connection between two entities
          does not by itself establish that they are
          necessarily related.

    \item \textbf{Identification does not establish
          ontological validity.}
          The existence of an identifier does not prove
          the existence, truth, or validity of the
          entity or claim to which it refers.

    \item \textbf{Registration does not establish
          philosophical acceptance.}
          Recording a claim in a project register does
          not imply that the claim has been accepted
          as true or adequately justified.
\end{itemize}

The reference system therefore organizes the documentation
of the philosophical investigation while leaving the
substantive evaluation of its concepts, commitments,
arguments, and conclusions to the relevant philosophical
analyses.


% =============================================================================
% 11 — RELATION TO OTHER PROJECT REFERENCES
% =============================================================================

\section*{Relation to Other Project References}

The Logical Entity Reference Map operates in coordination
with the other documentary and organizational components
of Book I.

Each component has a distinct function and must retain
its own responsibility within the project architecture.

\begin{center}
\small
\setlength{\tabcolsep}{4pt}
\begin{tabularx}{\textwidth}{
    >{\raggedright\arraybackslash}p{0.32\textwidth}
    >{\raggedright\arraybackslash}X
}
\toprule
\textbf{Project Component} &
\textbf{Relationship to This Map}
\\
\midrule

Acronyms and Structural Codes
&
Defines and explains the abbreviated forms and
structural codes used throughout the project.
\\

Glossary
&
Provides the terminology and adopted meanings
of terms used in the book.
\\

Asset Directories
&
Maintain the independent files associated with
the asset identifiers.
\\

Logical Registers
&
Maintain classified entries, their identifiers,
and their associated documentary records.
\\

Chapter and Section Files
&
Contain the substantive philosophical investigations
to which structural references point.
\\

Bibliography
&
Records the bibliographic sources cited or otherwise
included in the scholarly apparatus of the book.
\\

Indexes
&
Provide navigational access to selected names,
subjects, and concepts discussed in the text.
\\
\bottomrule
\end{tabularx}
\end{center}

The Logical Entity Reference Map coordinates these
components at the level of identification and reference.
It does not replace their respective functions.

In particular, the glossary remains responsible for
terminological clarification; the registers remain
responsible for maintaining classified records; the
asset directories remain responsible for independent
asset storage; and the chapters and sections remain
the primary locations of the philosophical arguments.

The map provides a common documentary framework through
which these distinct components can be located and
related without collapsing their respective roles.


% =============================================================================
% 12 — FINAL REFERENCE PRINCIPLE
% =============================================================================

\section*{Final Reference Principle}

The governing principle of the reference system is
the separation of identification from determination.

An identifier designates a documentary entity and
makes its location traceable within the project.
A classification specifies the functional category
assigned to an entry. A justificatory status records
the position assigned to that entry in the evaluation
of its philosophical standing.

None of these operations, considered merely as
documentary operations, establishes the ontological
constitution or philosophical validity of its referent.

The reference system must therefore remain precise
in its identifying function, consistent in its
classifications, and limited in the claims that
may legitimately be inferred from its structure.

Its purpose is to make the philosophical architecture
of Book I traceable and navigable without substituting
documentary organization for philosophical argument.


% =============================================================================
% END OF LOGICAL ENTITY REFERENCE MAP
% =============================================================================